% TOP_PROBLEM: {"id": 3226, "title": "Reed's omega, delta, and chi conjecture", "category_id": 3, "category": {"id": 3, "name": "graph_theory", "display_name": "Graph Theory", "description": "Problems involving graphs, networks, and their properties.", "slug": "graph-theory", "order_index": 3, "created_at": "2026-07-31T15:26:25.671Z"}, "status": "open", "problem_number": "OPG-335", "set_id": 12}
% FIRST_POSTED: 2026-10-01
% ENTRY_KIND: statement_only
% UnsolvedMath ID 3226; source catalogue code OPG-335
% !TeX program = lualatex
% Definition and sources only; this file is not an LLM solution attempt.
\documentclass[11pt,a4paper]{article}
\usepackage[margin=25mm]{geometry}
\usepackage{fontspec}
\setmainfont{lmroman10-regular.otf}[BoldFont=lmroman10-bold.otf,
 ItalicFont=lmroman10-italic.otf,BoldItalicFont=lmroman10-bolditalic.otf]
\usepackage{amsmath,amssymb,mathtools}
\usepackage{unicode-math}
\setmathfont{latinmodern-math.otf}
\AtBeginDocument{\let\setminus\smallsetminus}
\usepackage{xurl,hyperref}
\hypersetup{colorlinks=true,urlcolor=blue}
\setcounter{secnumdepth}{0}
\setlength{\parindent}{0pt}
\setlength{\parskip}{5pt}
\setlength{\emergencystretch}{3em}
\begin{document}
\section*{Reed's omega, delta, and chi conjecture}\label{3226}
{\small UnsolvedMath ID 3226 · OPG-335 · Graph Theory · OpenGarden}
\subsection{Definitions and mathematical statement}
Let $G=(V,E)$ be a nonempty
finite simple undirected graph.
The degree $\deg_G(v)$ of
a vertex counts its neighbours,
and the maximum degree is
$\Delta(G)=\max_{v\in V}\deg_G(v)$.
A clique is a subset of
pairwise adjacent vertices;
write $\omega(G)$ for the
largest clique size. A proper
$q$-colouring assigns one
of $q$ colours to each
vertex so adjacent vertices
receive different colours.
The chromatic number $\chi(G)$
is the least possible $q$.

Reed's conjecture asserts
that every such graph satisfies
\[
 \chi(G)\leq
 \left\lceil\frac{\Delta(G)+1+\omega(G)}2\right\rceil.
\]
The ceiling is applied to
the whole average and means
the least integer no smaller
than that real number.
All three parameters refer
to the same graph.

Every clique requires distinct
colours, so $\omega(G)\leq\chi(G)$.
Greedy colouring gives
$\chi(G)\leq\Delta(G)+1$.
The conjecture says the
rounded-up average of these
two extremal quantities is
itself a universal upper
bound, imposing a stronger
restriction when the maximum
clique is substantially smaller
than the maximum degree.

No connectedness, planarity,
or regularity is assumed.
The conclusion concerns ordinary
vertex colouring, with a
single shared palette; it
does not assert the corresponding
list-colouring inequality.
For an edgeless nonempty graph
the bound reads $1\leq1$.
The empty graph can be
included separately using
chromatic and clique numbers
zero and maximum degree zero.
\subsection{Short English statement}
Is every graph colourable with the rounded-up average of its clique number and one more than its maximum degree?
\subsection{Sources}
\begin{itemize}
\item[S1] ULAM.AI and the UnsolvedMath Contributors, supplied problems.json, record 3226 (OPG-335), OpenGarden collection. Catalogue identity and source transcription; CC BY 4.0.
\url{https://huggingface.co/datasets/ulamai/UnsolvedMath}.
\item[S2] Open Problem Garden, Reed's omega, delta, and chi conjecture. Original problem and discussion; the mathematical formulation below is expanded from the supplied source record.
\url{http://www.openproblemgarden.org/op/reeds_omega_delta_and_chi_conjecture}.
\end{itemize}
\end{document}
